\section{Digital zoom as an exact camera transform}
\label{sec:law}

\paragraph{The law.}
Let an image $I$ of size $W\times H$ be taken by a pinhole camera with intrinsics
$K=\begin{psmallmatrix} f_x & 0 & c_x\\ 0 & f_y & c_y\\ 0&0&1\end{psmallmatrix}$.
Cropping the box $b=(x_0,y_0,w,h)$ and resizing it back to $W\times H$ gives an image
$\mathcal{Z}_b I$ that is \emph{exactly} the image of the same scene taken with
\begin{equation}
  f'_x = s_x f_x,\quad c'_x = s_x\,(c_x - x_0 + \tfrac12) - \tfrac12,\qquad s_x = W/w,
  \label{eq:law}
\end{equation}
and likewise for $y$ (the half pixels come from pixel centres).
Every pixel that survives the crop sees the same 3D point, so its metric depth is unchanged.
For a centred crop with $s_x=s_y=s$ we call $s$ the \emph{zoom factor}.
\cref{eq:law} needs no assumption about the scene, the model or the data;
our implementation reproduces it to below $10^{-3}$\,px (supplement).

\paragraph{What a correct model must do.}
Let a model map an image to a depth map $\hat D(I)$ and, optionally, a focal estimate
$\hat f(I)$. Zoom equivariance asks
\begin{equation}
  \hat D(\mathcal{Z}_b I)\big(\mathcal{Z}_b(p)\big) = \hat D(I)(p), \qquad
  \hat f(\mathcal{Z}_b I) = s\,\hat f(I)
  \label{eq:equiv}
\end{equation}
for all pixels $p$ inside $b$. Neither side involves ground truth.

\paragraph{Two exponents.}
Real models do not satisfy \cref{eq:equiv}; over $s\in[1,3]$ we summarise their behaviour by two power laws,
\begin{equation}
  \frac{\hat D(\mathcal{Z}_b I)}{\hat D(I)} \approx s^{-\alpha}, \qquad
  \frac{\hat f(\mathcal{Z}_b I)}{\hat f(I)} \approx s^{\beta},
  \label{eq:exponents}
\end{equation}
fitted per image by least squares in log-log space through the origin
(the depth ratio is the median over the pixels in $b$).
An equivariant model has $\alpha=0$ and $\beta=1$. The fit is an approximation: the median residual is
\fitRmsMedian{} in log depth, and $\alpha$ grows with the zoom (Depth Pro \DepthproAlphaSmall{} for $s\le1.5$,
\DepthproAlphaLarge{} for $s\ge2$; supplement).
A model that ignores the zoom in its camera estimate has $\beta=0$.

\paragraph{Three invariants.}
Unprojecting the prediction with the model's own focal gives a point cloud.
Under a zoom, depth scales by $s^{-\alpha}$ and the lateral coordinate
$X=(u-c)\,\hat D/\hat f$ by $s\cdot s^{-\alpha}/s^{\beta}=s^{1-\alpha-\beta}$. Hence
\begin{itemize}
  \item $\alpha=0$: depth is preserved;
  \item $\beta=1$: the shape of the point cloud (the ratio $X/Z$) is preserved;
  \item $\alpha+\beta=1$: lateral coordinates are preserved (not depth, not shape).
\end{itemize}
The third relation is what a model obtains if it reasons from apparent size.
An object of metric width $L$ that spans $\ell$ pixels is placed at
$\hat D=\hat f L/\ell$. A zoom multiplies $\ell$ by $s$ and $\hat f$ by $s^{\beta}$, so
$\hat D$ scales by $s^{\beta-1}$, that is $\alpha=1-\beta$.
A pixels-only model ($\hat f$ fixed, implicitly) that reasons from size has $\alpha=1$:
twice as large means twice as close.

\paragraph{The zoom--squash ambiguity.}
Why would a model reason from size at all, when the zoom is in plain sight?
Because a zoom is invisible to size cues.
A point $(X,Y,Z)$ seen with focal $sf$ projects to $u = sf\,X/Z + c = f\,X/(Z/s) + c$.
\begin{quote}
\emph{A zoomed photo of a scene is pixel-for-pixel the unzoomed photo of the same scene
squashed along the viewing axis by $s$.}
\end{quote}
This is the classical focal--depth ambiguity of a single view~\cite{he2018focal}, written as an exact, testable
statement about a zoom; it is not a new ambiguity.
The squashed scene has the same lateral sizes, so every object has a plausible width; only its
extent \emph{in depth} is wrong (a 4.5\,m car becomes 2.25\,m long at $s=2$), and so are the
angles between surfaces. Telling the two apart needs priors on 3D shape (aspect ratios in
depth, right angles, the rate at which parallel lines converge), not priors on size.
One consequence is a hard limit: a fronto-parallel plane squashed along the viewing axis is still a fronto-parallel plane, so
\emph{no} estimator can be equivariant on such scenes without being told the camera.
The ambiguity also explains why the third invariant is the natural failure mode: believing the
squashed scene means $\alpha=1$ with $\beta=0$, a point on the line $\alpha+\beta=1$.
Note that the line only guarantees lateral consistency; along it, depth and shape are wrong
unless $\beta=1$.

\paragraph{Models that take the camera as input.}
Models such as Metric3D~\cite{yin2023metric3d,hu2024metric3dv2} receive $K$ and predict depth in a
canonical camera. With the correct $K'$ from \cref{eq:law} they are equivariant by design.
In practice the zoom often happens after capture (a crop in a gallery app, a screenshot,
a detection crop) and the metadata still describes the original camera~\cite{kim2026mismatch}.
In this \emph{stale-$K$} setting the network sees the same pixels and only the final
de-canonical scale $f/1000$ differs by the factor $s$, so $\alpha_{\text{stale}}=\alpha_{\text{true}}+1$
exactly, up to the model's output clipping. Our runs confirm this (mean difference over \staleImages{} images
\staleDiffMean), so a camera-conditioned model with stale metadata behaves like a pixels-only model that
reads apparent size.
